\documentclass[11pt,a4paper]{article}

% -----------------------------------------------------------------------------
% Packages
% -----------------------------------------------------------------------------
\usepackage[margin=1in]{geometry}
\usepackage{amsmath,amssymb,amsfonts}
\usepackage{bm}
\usepackage{booktabs}
\usepackage{hyperref}
\usepackage{listings}
\usepackage{xcolor}
\usepackage{enumitem}
\usepackage{graphicx}
\usepackage{microtype}
\usepackage{tikz}
\usetikzlibrary{arrows.meta,positioning,fit,calc,shapes.geometric}

\hypersetup{
  colorlinks=true,
  linkcolor=blue!60!black,
  urlcolor=blue!60!black,
  citecolor=blue!60!black
}

\lstset{
  basicstyle=\ttfamily\small,
  breaklines=true,
  frame=single,
  backgroundcolor=\color{gray!8},
  keywordstyle=\color{blue!60!black}\bfseries,
  stringstyle=\color{teal!70!black},
  commentstyle=\color{gray},
  showstringspaces=false,
  columns=fullflexible,
  keepspaces=true
}

\title{Methodology of the PlasMol Fourier Driver\\
\large Including Meep Vacuum Field References for Hybrid Spectra}
\author{PlasMol Development Notes}
\date{\today}

\begin{document}
\maketitle

\tableofcontents
\newpage

% =============================================================================
\section{Purpose and scope}
% =============================================================================

The \textbf{Fourier driver} is the recommended workflow for computing linear
absorption spectra.  Given a short, broadband excitation, the code:
\begin{enumerate}
  \item runs three independent polarizations ($x$, $y$, $z$) in parallel;
  \item records the induced molecular dipole $\boldsymbol{\mu}(t)$;
  \item optionally damps the time-domain signals to artificially increase the likelihood of convergence;
  \item Fourier-transforms the response;
  \item assembles a normalized absorption spectrum over a user-chosen energy
        window.
\end{enumerate}

% =============================================================================
\section{Physical overview}
% =============================================================================

\subsection{Induced dipole from RT-TDDFT}

At each quantum step the molecule is propagated under the instantaneous
electric field $\mathbf{E}(t)$ (either an analytic kick/pulse or the local
Meep field).  The induced dipole is
\begin{equation}
  \mu_i(t)
  =
  \operatorname{Tr}\!\big[
    \hat{\mu}_i \big(\mathbf{D}(t) - \mathbf{D}_0\big)
  \big],
  \qquad i \in \{x,y,z\},
  \label{eq:mu}
\end{equation}
where $\mathbf{D}(t)$ is the time-dependent density matrix in the AO basis, $\mathbf{D}_0$ is
the ground-state density, and $\hat{\mu}_i$ are dipole integrals
(\texttt{molecule.calculate\_mu()}).  Open-shell runs sum $\alpha$ and $\beta$
blocks before the trace.

\subsection{From time series to absorption}

PlasMol builds a scalar absorption-like quantity from the imaginary part of that response
along all three Cartesian directions and a prefactor proportional to frequency:
\begin{equation}
  A(\omega)
  =
  -\frac{4\pi\omega}{3c}
  \sum_{i\in\{x,y,z\}}
  \operatorname{Im}\!\big[\mathcal{R}_i(\omega)\big],
  \label{eq:abs}
\end{equation}
where $c$ is the speed of light in atomic units
(\texttt{constants.C\_AU}) and $\mathcal{R}_i(\omega)$ is either the
Fourier-transformed dipole or the deconvolved ratio
$\mu_i(\omega)/E_{\mathrm{inc},ii}(\omega)$ depending on the source kick; see
Sections~\ref{sec:qm}--\ref{sec:deconv}.  The published curve is
peak-normalized:
\begin{equation}
  A_{\mathrm{norm}}(\omega)
  =
  \frac{A(\omega)}{\max_\omega |A(\omega)|}.
\end{equation}

\noindent Frequencies are obtained from the real FFT frequency grid,
\begin{equation}
  \omega_k = 2\pi\, f_k,
  \qquad
  f_k = \mathrm{fftfreq}(N,\Delta t),
\end{equation}
converted to eV via $E[\mathrm{eV}] = \omega[\mathrm{a.u.}] \times 27.211386$,
and restricted to the user window
$[\texttt{min\_ev},\,\texttt{max\_ev}]$.

% =============================================================================
\section{Common post-processing pipeline}
\label{sec:common}
% =============================================================================

Regardless of how the directional runs are obtained (Quantum only or Full Plasmol), the post-processing
stage is shared.

\subsection{Folding diagonal response components}

Each directional run writes a full Cartesian field CSV.  Only the component
parallel to the drive is retained and stacked into a $(3\times N)$ array:
\begin{align}
  \boldsymbol{\mu}_{\mathrm{diag}}(t)
  &=
  \big(
    \mu_x^{(x)}(t),
    \mu_y^{(y)}(t),
    \mu_z^{(z)}(t)
  \big)^\top
  \qquad\text{(function \texttt{fold})}.
\end{align}
Here the superscript denotes the polarization of the incident source.  Time
grids are required to share a common $\Delta t$; series are trimmed to the
shortest length if they differ slightly.

\subsection{Optional exponential damping}

Two related damping mechanisms may be applied before the FFT.

\paragraph{$\gamma$ window.}
Inside \texttt{fourier()}, every time series is multiplied by
\begin{equation}
  w_\gamma(t) = e^{-\gamma t},
\end{equation}
with $\gamma = \texttt{fourier.gamma}$ (a.u.; default = 0). This broadens spectral
features artificially.

\paragraph{$\tau$ damping.}
If $\texttt{fourier.tau} > 0$, an additional factor
\begin{equation}
  w_\tau(t) = e^{-t/\tau}
\end{equation}
is applied to the polarizability CSVs on disk (written as
\texttt{*\_damped.csv}) and, for the Meep path, to the reference
$\mathbf{E}_{\mathrm{inc}}$ in memory only so that the reusable
reference file remains undamped.

\subsection{Discrete Fourier transform}

For each Cartesian axis $i$,
\begin{equation}
  S_{\mu,i}(\omega)
  =
  \Delta t \sum_{n=0}^{N-1}
  \mu_i(t_n)\, w(t_n)\, e^{-i\omega t_n},
  \label{eq:fft}
\end{equation}
implemented as \texttt{np.fft.fft(...) * dt}.  The same window $w$ multiplies
$\mathbf{E}_{\mathrm{inc}}$ when deconvolution is active.

\subsection{Outputs}

\begin{itemize}
  \item \texttt{spectrum\_filepath} --- PNG of $A_{\mathrm{norm}}(E)$;
  \item sibling \texttt{.csv} with columns \texttt{Frequency}, \texttt{Absorption};
  \item optional \texttt{npz\_filepath} storing
        \texttt{abs\_imag}, \texttt{abs\_real}, \texttt{freqs}, and a
        \texttt{deconvolved} flag;
  \item per-direction working directories \texttt{x\_dir/}, \texttt{y\_dir/},
        \texttt{z\_dir/} with raw \texttt{field\_e.csv} /
        \texttt{field\_p.csv} (and vacuum reference CSVs when applicable).
\end{itemize}

% =============================================================================
\section{Quantum-only Fourier path (delta kicks)}
\label{sec:qm}
% =============================================================================

\subsection{Setup}

When the input has a \texttt{molecule} section but no \texttt{plasmon}
section, the Fourier driver schedules three pure RT-TDDFT jobs.
For each direction $i\in\{x,y,z\}$:
\begin{itemize}
  \item \texttt{molecule.source.component} is forced to $i$;
  \item \texttt{QUANTUMSOURCE} builds an analytic field
        $\mathbf{E}^{(i)}(t)$ on the global time grid
        $t_n = n\,\Delta t$;
  \item for \texttt{type:"kick"}, the active component is a one-step
        $\delta$-like impulse of strength
        \texttt{intensity} at \texttt{peak\_time}.
\end{itemize}

\noindent A true $\delta$-kick has a flat spectrum, so the polarizability is
proportional to the Fourier transform of the induced dipole itself.  The
driver therefore calls
\begin{equation}
  \mathcal{R}_i(\omega) = S_{\mu,i}(\omega)
  \qquad\text{(\texttt{field\_e = None} branch of \texttt{fourier()})},
\end{equation}
i.e.\ $\operatorname{Im}[\mu_i(\omega)]$ enters Eq.~\eqref{eq:abs} directly.

\subsection{Typical JSON fragment}

\begin{lstlisting}[language={}]
{
  "settings": {
    "dt": 0.1,
    "t_end": 4000,
    "driver": "fourier"
  },
  "molecule": {
    "geometry": [
        {"atom": "O", "coord": [ 0.000, 0.000, -0.130]},
        {"atom": "H", "coord": [ 1.489, 0.000,  1.033]},
        {"atom": "H", "coord": [-1.489, 0.000,  1.033]}
    ],
    "geometry_units": "bohr",
    "basis": "631g*",
    "xc": "pbe0",
    "source": {
      "type": "kick",
      "intensity": 5e-5,
      "peak_time": 0.05,
      "width_steps": 3
    }
  },
  "files": {
    "checkpoint": {
      "filepath": "checkpoint.npz",
      "frequency_steps": 1000
    },
    "field_e_filepath": "field_e.csv",
    "field_p_filepath": "field_p.csv",
    "spectra_e_vs_p_filepath": "output.png"
  },
  "additional_parameters": {
    "fourier": {
      "min_ev": 1,
      "max_ev": 35,
      "npz_filepath": "fourier.npz",
      "spectrum_filepath": "spectrum.png"
    }
  }
}
\end{lstlisting}

% =============================================================================
\section{Hybrid Meep Fourier path}
\label{sec:hybrid}
% =============================================================================

\subsection{Motivation}

When a nanoparticle (or any classical structure) is present, the molecule no
longer sits in a free-space $\delta$-kick.  Instead Meep propagates a broadband
source (typically a \texttt{gaussian} pulse) through the FDTD grid; the local
field at \texttt{plasmon.molecule.position} drives RT-TDDFT, and the induced
dipole may be fed back into Meep as a point \texttt{CustomSource}
(\texttt{back\_propagation}).  Two complications follow:
\begin{enumerate}
  \item the incident pulse has a non-flat spectrum and a non-trivial phase;
  \item the local field at the molecule includes scattering from the
        nanoparticle, which is precisely the physics one often wants to
        keep in $\boldsymbol{\mu}(t)$.
\end{enumerate}
Dividing $\mu(\omega)$ by the \emph{total} local field would remove
plasmon-enhanced response.  Dividing by a pure \emph{vacuum incident} field
$\mathbf{E}_{\mathrm{inc}}$ removes only the source envelope, leaving
nanoparticle / local-field physics inside $\boldsymbol{\mu}$.

\subsection{Production runs (molecule + NP)}

For each polarization $i$, the driver sets \texttt{plasmon.source.component = i}, 
and writes outputs under \texttt{x\_dir/}, \texttt{y\_dir/}, or \texttt{z\_dir/}.  

\noindent At every Meep timestep (and at $t=0$):
\begin{enumerate}
  \item sample $\mathbf{E}$ at the molecule site
        (\texttt{\_get\_electric\_field}), converting Meep units to atomic
        units via \texttt{convertFieldMeep2Atomic};
  \item if $|\mathbf{E}|$ exceeds \texttt{tolerance\_field\_e}, call
        RT-TDDFT (\texttt{propagation}) to obtain $\boldsymbol{\mu}$;
  \item store $\boldsymbol{\mu}$ for the next Meep source query (if
        back-coupling is on) and append both $\mathbf{E}$ and
        $\boldsymbol{\mu}$ to CSV.
\end{enumerate}

\noindent Additionally, an identical simulation begins but without the NP nor the molecule 
to calculate the reference E field spectra.
\begin{enumerate}
  \item evaluates $\mathbf{E}(\mathbf{r}_{\mathrm{mol}}, t)$ in atomic units;
  \item appends $(t, E_x, E_y, E_z)$ to the configured CSV.
  \item once all three runs are done, the diagonal elements are merged into \( \mathbf{E}_{\text{diag}}(t) \).
\end{enumerate}

\subsection{Unit and time-grid notes}

Hybrid timesteps are conformed to Meep's Courant-limited $\Delta t$:
\begin{equation}
  \Delta t_{\mathrm{Meep}} = \frac{\mathrm{Courant}}{\mathrm{resolution}},
  \qquad
  \Delta t_{\mathrm{a.u.}}
  =
  \Delta t_{\mathrm{Meep}}
  \times
  \texttt{convertTimeMeep2Atomic}.
\end{equation}

\subsection{Spectral deconvolution}
\label{sec:deconv}

This subsection develops the hybrid Fourier formula from linear response,
states precisely what is divided by what, and connects the continuous
theory to the discrete operations in \texttt{fourier()}.

Solving for the polarizability is therefore a complex division,
\begin{equation}
  \alpha_i(\omega)
  =
  \frac{\mu_i(\omega)}{E_i^{\mathrm{drive}}(\omega)},
  \label{eq:alpha-def}
\end{equation}
whenever the denominator is non-vanishing.  The absorption spectrum is
built from $\operatorname{Im}\alpha_i(\omega)$ (Section~\ref{sec:common}
and Eq.~\eqref{eq:abs}), not from $\operatorname{Im}\mu_i(\omega)$ alone:
the latter still carries the spectral envelope of the Meep source.

If $E_i^{\mathrm{drive}}(t) = \kappa\,\delta(t)$, then
$E_i^{\mathrm{drive}}(\omega) = \kappa$ (flat in frequency), so
\begin{equation}
  \alpha_i(\omega)
  =
  \frac{\mu_i(\omega)}{\kappa}
  \propto
  \mu_i(\omega).
\end{equation}
Hence the pure quantum path may take
$\mathcal{R}_i(\omega) = S_{\mu,i}(\omega)$ up to a global scale that
disappears after peak normalization.  A Meep Gaussian (or other finite
pulse) has $E_i^{\mathrm{drive}}(\omega)$ that is neither flat nor
real; omitting the division imprints the pulse spectrum and phase onto
$A(\omega)$.

Write the electric field at the molecule site in the Full Plasmol run as
\begin{equation}
  \mathbf{E}_{\mathrm{loc}}(\mathbf{r}_{\mathrm{mol}}, t)
  =
  \mathbf{E}_{\mathrm{inc}}(t)
  +
  \mathbf{E}_{\mathrm{scat}}(t),
  \label{eq:E-split}
\end{equation}
where $\mathbf{E}_{\mathrm{inc}}$ is the free-space incident field of the
Meep source (same cell, source, and sample point, but no nanoparticle and
no quantum molecule) and $\mathbf{E}_{\mathrm{scat}}$ collects everything
scattered by the nanoparticle (and, if back-coupling is on, radiation
from the molecular dipole itself in later times).\\

\noindent Two natural choices for $E_i^{\mathrm{drive}}$ in Eq.~\eqref{eq:alpha-def}
have very different meanings:

\begin{center}
\begin{tabular}{@{}lp{8.6cm}@{}}
\toprule
Denominator & What $\alpha_i$ measures \\
\midrule
$E_{\mathrm{loc},i}(\omega)$
  & Response relative to the \emph{total} local field.  Plasmonic
    near-field enhancement largely cancels between numerator and
    denominator; peaks track bare molecular resonances. \\
$E_{\mathrm{inc},ii}(\omega)$
  & Response relative to the \emph{vacuum incident} field.  Enhancement
    and scattering remain inside $\mu_i(\omega)$; $A(\omega)$ reports
    the molecule as driven by the same free-space pulse that would be
    present without the NP. \\
\bottomrule
\end{tabular}
\end{center}

\noindent PlasMol uses the second choice.  Operationally, three vacuum Meep runs
record $\mathbf{E}_{\mathrm{inc}}^{(i)}(t)$ for $i\in\{x,y,z\}$; only the
co-polarized diagonal is retained:
\begin{equation}
  \mathbf{E}_{\mathrm{diag}}(t)
  =
  \begin{pmatrix}
    E_x^{(x)}(t) \\
    E_y^{(y)}(t) \\
    E_z^{(z)}(t)
  \end{pmatrix}
  \equiv
  \begin{pmatrix}
    E_{\mathrm{inc},xx}(t) \\
    E_{\mathrm{inc},yy}(t) \\
    E_{\mathrm{inc},zz}(t)
  \end{pmatrix}.
  \label{eq:E-diag}
\end{equation}
The production runs contribute the analogous diagonal dipole
$\boldsymbol{\mu}_{\mathrm{diag}}(t)$ from \texttt{fold}.

\subsubsection{Consistent deconvolution}

Define the discrete Fourier transforms used in code
($t_n = t_0 + n\Delta t$, $n=0,\ldots,N-1$):
\begin{align}
  S_{\mu,i}(\omega_k)
  &=
  \Delta t
  \sum_{n=0}^{N-1}
  \mu_i(t_n)\, w(t_n)\,
  e^{-2\pi i\, kn/N},
  \label{eq:Smu} \\
  S_{E,i}(\omega_k)
  &=
  \Delta t
  \sum_{n=0}^{N-1}
  E_{\mathrm{inc},ii}(t_n)\, w(t_n)\,
  e^{-2\pi i\, kn/N},
  \label{eq:SE}
\end{align}
where $\omega_k = 2\pi\, \mathrm{fftfreq}(N,\Delta t)_k$ in atomic units.
Because the \emph{same} window multiplies both time series, it factors in
a manner compatible with Eq.~\eqref{eq:alpha-def}: if
$\mu = \chi * E_{\mathrm{inc}}$ and one multiplies both signals by $w$
before transforming, the ratio
\begin{equation}
  \hat\alpha_i(\omega_k)
  :=
  \frac{S_{\mu,i}(\omega_k)}{S_{E,i}(\omega_k)}
  \label{eq:alpha-hat}
\end{equation}
still estimates $\alpha_i(\omega_k)$ up to the usual finite-record and
window-bias errors of a truncated FFT---it does \emph{not} introduce an
extra $w(\omega)$ factor in the numerator alone.  (If only $\mu$ were
windowed, the ratio would mix the polarizability with the Fourier
transform of the window.)

\subsubsection{Complex division, phase cancellation, and the amplitude floor}

Write the complex spectra in polar form,
\begin{equation}
  S_{\mu,i}(\omega)
  =
  \big|S_{\mu,i}\big|\, e^{i\phi_{\mu,i}},
  \qquad
  S_{E,i}(\omega)
  =
  \big|S_{E,i}\big|\, e^{i\phi_{E,i}}.
\end{equation}
Then
\begin{equation}
  \hat\alpha_i(\omega)
  =
  \frac{\big|S_{\mu,i}\big|}{\big|S_{E,i}\big|}
  \,
  e^{i(\phi_{\mu,i}-\phi_{E,i})}.
  \label{eq:polar}
\end{equation}
Two useful consequences follow immediately:
\begin{enumerate}
  \item \textbf{Amplitude flattening.}
        Division by $|S_{E,i}|$ removes the Meep source envelope
        (e.g.\ the Gaussian spectrum set by \texttt{frequency} /
        \texttt{fwidth}).  Residual structure in
        $|\hat\alpha_i|$ is molecular and/or plasmon-mediated.
  \item \textbf{Phase cancellation.}
        The relative phase $\phi_{\mu,i}-\phi_{E,i}$ is the phase lag of
        the induced dipole relative to the vacuum incident field.
        Propagation delay from the Meep source plane to
        $\mathbf{r}_{\mathrm{mol}}$, and the overall complex scale of the
        source, largely cancel.  Absorption depends on
        $\operatorname{Im}\hat\alpha_i$, which is the component of the
        response in phase quadrature with $E_{\mathrm{inc}}$.
\end{enumerate}

In components,
\begin{align}
  \operatorname{Re}\hat\alpha_i
  &=
  \frac{
    \operatorname{Re}S_{\mu,i}\,\operatorname{Re}S_{E,i}
    +
    \operatorname{Im}S_{\mu,i}\,\operatorname{Im}S_{E,i}
  }{
    |S_{E,i}|^2
  },
  \label{eq:re-alpha} \\[0.4em]
  \operatorname{Im}\hat\alpha_i
  &=
  \frac{
    \operatorname{Im}S_{\mu,i}\,\operatorname{Re}S_{E,i}
    -
    \operatorname{Re}S_{\mu,i}\,\operatorname{Im}S_{E,i}
  }{
    |S_{E,i}|^2
  }.
  \label{eq:im-alpha}
\end{align}
Equation~\eqref{eq:im-alpha} is the quantity summed into $A(\omega)$.
The code stores both real and imaginary parts (\texttt{abs\_real},
\texttt{abs\_imag} in the optional NPZ); only the imaginary part enters
\texttt{absorption()}.

Near frequencies where the vacuum pulse has little power,
$|S_{E,i}|\to 0$ and Eq.~\eqref{eq:alpha-hat} is ill-conditioned.
PlasMol therefore applies a relative amplitude floor on the spectral
window of interest only.  Let
\begin{equation}
  \mathcal{W}
  =
  \big\{\, k :\;
    \texttt{min\_ev}
    \le
    \hbar\omega_k
    \le
    \texttt{max\_ev}
  \,\big\}
\end{equation}
(with $\hbar\omega$ converted to eV as in Section~\ref{sec:common}), and
\begin{equation}
  E_{\max,i}
  =
  \max_{k\in\mathcal{W}}
  \big|S_{E,i}(\omega_k)\big|,
  \qquad
  \varepsilon
  =
  \varepsilon_{\mathrm{rel}}\, E_{\max,i},
  \qquad
  \varepsilon_{\mathrm{rel}} = 10^{-8}
  \;\;(\texttt{e\_floor\_rel}).
\end{equation}
The regularized estimator is
\begin{equation}
  \alpha_i(\omega_k)
  =
  \begin{cases}
    \displaystyle
    \dfrac{S_{\mu,i}(\omega_k)}{S_{E,i}(\omega_k)}
      & \text{if } k\in\mathcal{W}
           \text{ and }
           |S_{E,i}(\omega_k)| > \varepsilon, \\[0.8em]
    0
      & \text{otherwise}.
  \end{cases}
  \label{eq:alpha}
\end{equation}
If no bin in $\mathcal{W}$ exceeds the floor for a given polarization,
that direction contributes nothing and a warning is logged.  In practice
this means: either widen the Meep source bandwidth so that
$|E_{\mathrm{inc}}(\omega)|$ covers $[\texttt{min\_ev},\texttt{max\_ev}]$,
or shrink the plot window to the support of the pulse.

\subsubsection{From \texorpdfstring{$\alpha_i$}{alpha} to the reported absorption}

With $\mathcal{R}_i(\omega) := \alpha_i(\omega)$ in the hybrid path,
Eq.~\eqref{eq:abs} becomes
\begin{equation}
  A(\omega)
  =
  -\frac{4\pi\omega}{3c}
  \sum_{i\in\{x,y,z\}}
  \operatorname{Im}\alpha_i(\omega).
  \label{eq:abs-hybrid}
\end{equation}
For an isotropic molecule in free space one recovers the familiar
orientational average of the dipole polarizability,
$\tfrac13\sum_i\operatorname{Im}\alpha_{ii}$.  With a nanoparticle the
same algebraic form is retained, but each $\alpha_i$ is an
\emph{effective} polarizability of the molecule as driven by the vacuum
incident field in that geometry---not the bare gas-phase $\alpha$.

Finally the curve is peak-normalized for shape comparison,
\begin{equation}
  A_{\mathrm{norm}}(\omega)
  =
  \frac{A(\omega)}{\max_{\omega\in\mathcal{W}} |A(\omega)|},
\end{equation}
so overall scales (kick strength, Meep amplitude, unit conventions
shared by $\mu$ and $E_{\mathrm{inc}}$) drop out, provided they are
consistent between numerator and denominator.

\subsubsection{Compact algorithm (as implemented)}

\begin{enumerate}
  \item Align $\boldsymbol{\mu}_{\mathrm{diag}}(t)$ and
        $\mathbf{E}_{\mathrm{diag}}(t)$ on a common length and $\Delta t$
        (\texttt{\_align\_series}).
  \item Form $w(t)$ from $\gamma$ (and $\tau$ if set).
  \item For each $i\in\{x,y,z\}$:
        compute $S_{\mu,i}$ and $S_{E,i}$ by Eqs.~\eqref{eq:Smu}--\eqref{eq:SE};
        restrict to $\mathcal{W}$; apply Eq.~\eqref{eq:alpha}.
  \item Assemble $A(\omega)$ from Eq.~\eqref{eq:abs-hybrid}; normalize;
        write PNG/CSV and optional NPZ (\texttt{deconvolved=True}).
\end{enumerate}

\paragraph{Summary of interpretation.}
\begin{itemize}
  \item Division by vacuum $E_{\mathrm{inc}}(\omega)$ flattens the Meep
        source spectrum and cancels its phase
        (Eqs.~\eqref{eq:polar}--\eqref{eq:im-alpha}).
  \item Nanoparticle-induced local fields remain inside $\mu(\omega)$, so
        plasmon-enhanced features are retained
        (Eq.~\eqref{eq:E-split} and the table above).
  \item Bins with negligible $|E_{\mathrm{inc}}|$ are zeroed rather than
        producing numerical blow-ups (Eq.~\eqref{eq:alpha}).
\end{itemize}

\subsection{Example hybrid JSON (Fourier + Au NP)}

\begin{lstlisting}[language={}]
{
  "settings": {
    "dt": 0.1,
    "t_end": 10000,
    "driver": "fourier"
  },
  "plasmon": {
    "simulation": {
      "cell_length": 0.2,
      "pml_thickness": 0.05
    },
    "source": {
      "type": "gaussian",
      "center": [-0.04, 0, 0],
      "size": [0, 0.2, 0.2],
      "component": "z",
      "is_integrated": true,
      "additional_parameters": {
        "frequency": 2.291666667,
        "fwidth": 2.083333335
      }
    },
    "nanoparticle": {
      "material": "Au",
      "radius": 0.025
    },
    "molecule": {
      "position": [0, 0, 0.026451]
    }
  },
  "molecule": {
    "geometry": "Na.xyz",
    "geometry_units": "bohr",
    "charge": 0,
    "spin": 1,
    "basis": "631g*",
    "xc": "pbe0",
    "hermiticity_tolerance": 1e-12,
  },
  "files": {
    "field_e_filepath": "field_e.csv",
    "field_p_filepath": "field_p.csv",
    "spectra_e_vs_p_filepath": "output.png"
  },
  "additional_parameters": {
    "fourier": {
      "spectrum_filepath": "spectrum.png",
      "min_ev": 1.5,
      "max_ev": 5.0,
      "field_e_ref_filepath": "field_e_ref.csv"
    }
  }
}
\end{lstlisting}

\end{document}